\chapter{石油、天然气与燃料}

\begin{kaichang}
傍晚的厨房里有两团火。灶台上，燃气灶的火苗安静地烧着，是干净的蓝色，几乎不发光，只把空气烤得微微颤动；桌角过生日剩下的蜡烛，火苗却是明亮的黄色，跳来跳去，顶端还拖着一缕几乎看不见的黑烟。

再看一眼打火机：透明外壳里装着半罐无色液体，晃一晃像水，按下去喷出来的却是能点燃的气。

现在请你猜三个问题：蓝色的火苗和黄色的火苗，烧的东西有什么不一样？打火机里的液体，怎么一出来就变成了气？还有——据说天然气本身一点味道都没有，可家里燃气泄漏时却能闻到刺鼻的臭味，这臭味是哪儿来的？

把猜测写下来。这一章结束时，这三个问题会串成同一条答案。
\end{kaichang}

\section{先看清火焰}

先不急着钻进分子，只用眼睛、鼻子和一只冷盘子，看看燃烧这件事本身。

把一只冰冷的瓷盘悬在燃气灶火苗上方十几厘米处（小心烫），盘底很快蒙上一层水珠——火焰在制造水。把同一只盘子凑近蜡烛火苗的上方，盘底除了水珠，还会蹭上一层黑灰，用手指一抹，一道黑印。蓝色的火几乎不留黑灰，黄色的火留。

再观察火苗的形状。蜡烛火焰的底部有一小圈暗蓝色，往上是明亮的黄色主体，最顶端颜色发暗。燃气灶的火则是整齐的蓝锥。如果你把燃气灶的风门调小，让空气进得不那么足，蓝火苗会慢慢变黄、变长，锅底也开始积黑——老式煤气灶熏黑锅底，就是这么来的。

鼻子也能提供线索。纯的汽油、酒精烧起来几乎闻不到烟味；蜡烛吹灭的一瞬间，那缕白烟有一股呛人的味道；柴油机的尾气更是又黑又呛。同样叫“燃烧”，产物的差别全写在颜色、气味和黑灰里。

\begin{shiyan}
\textbf{给火焰“拓一张指纹”}

材料：一根蜡烛、一个金属勺子（或瓷碟）、火柴、一杯水、成人陪同。

步骤：
\begin{enumerate}[itemsep=1pt]
\item 请成人点燃蜡烛，放在远离窗帘、纸张的桌面上。
\item 手握勺柄，把勺背伸进火焰外焰的上部，停留两三秒立即移开。
\item 观察勺背：先出现一层水雾，擦干后能看到一层黑色粉末，用手指捻一捻。
\item 再把勺子移到火焰最上方、离焰心较远处重复一次，比较两次黑粉的多少。
\end{enumerate}

安全提示：只烧两三秒，勺背会烫，不要用手碰；做完立刻吹灭蜡烛；全程需要成人在场。这层黑粉几乎就是纯碳——它是蜡烛“没烧完”的证据，记住它，本章后面还要跟它算账。
\end{shiyan}

\section{燃料的真身：碳和氢的链条}

这些火焰里烧的究竟是什么？把镜头推进到分子层，答案出人意料地单调：天然气、打火机气、汽油、煤油、柴油、蜡烛、沥青——这一大家子的主要成分，都是\textbf{烷烃}。

第 36 章讲过碳的四只“手”：每个碳原子能伸出 4 个单键，彼此连成链、连成环。烷烃是最简单的一类：碳与碳之间只用单键相连，剩下的“手”全部抓满氢原子。一个碳配四个氢，是甲烷 \chem{CH_4}，天然气的主角；两个碳连成小链，配六个氢，是乙烷 \chem{C_2H_6}；三个碳是丙烷，四个碳是丁烷——打火机里那罐“像水的液体”，就是加压后液化的丁烷 \chem{C_4H_{10}}。碳链每加长一节，就添上一个 \chem{CH_2}，像乐高一样无限延伸：戊烷、己烷、庚烷……一直排到几十个碳的长链。

因为碳的“手”全被占满，没有空余，烷烃被叫做\textbf{饱和烃}。占得满，意味着反应接口少（第 37 章说过，官能团才是分子的反应接口，烷烃几乎没有接口），所以烷烃性格温和：不跟酸碱随便反应，不容易变色变质，能安安稳稳在油罐里存上好几年。它一生只热衷干一件事——燃烧。这恰好是燃料最想要的性格：平时老实，点火才发力。

烷烃还有一对决定命运的性质：\textbf{分子极性极小，分子间只靠微弱的范德华力抱团}（第 22 章讲过这种分子间的“弱握手”）。后果有两个。第一，它们不溶于水——油和水永远分层，加油站地面上的油花雨水冲不进下水道，只会漂着走。第二，沸点几乎完全由链长决定：链越短，分子间抓得越松，越容易变成气体。

\begin{table}[h]
\centering
\begin{tabular}{clll}
\toprule
碳原子数 & 常温状态 & 沸点范围 & 日常用途 \\
\midrule
1～4 & 气体 & 低于 $0\,^{\circ}\mathrm{C}$ 左右 & 天然气、打火机气、液化石油气 \\
5～12 & 易挥发液体 & 约 $30$～$200\,^{\circ}\mathrm{C}$ & 汽油、溶剂 \\
10～16 & 液体 & 约 $150$～$300\,^{\circ}\mathrm{C}$ & 煤油、航空燃料 \\
14～20 & 黏稠液体 & 约 $250$～$350\,^{\circ}\mathrm{C}$ & 柴油 \\
20 以上 & 半固体到固体 & 更高 & 润滑油、石蜡、沥青 \\
\bottomrule
\end{tabular}
\caption{烷烃家族按链长排班。沸点范围是大致的，因为里面还混着支链异构体。}
\end{table}

这张表就是整章的地图：天然气是链最短的一截，汽油是中间的一截，蜡烛和沥青是最长的一截。它们不是不同的“原料”，是同一家族按沸点排开的队伍。

顺便纠正一个直觉：石油并不只是“烧掉的”。你今天身上至少有十几样东西直接来自它——化纤衣服、塑料水杯、运动鞋底、洗发水瓶、药膏、合成橡胶轮胎、铺路沥青。对现代工业来说，石油更宝贵的身份是\textbf{化工原料}：把它的碳链拆开、重组，就能拼出塑料、纤维、药物和染料的骨架（第 41 章会讲染料，第 48 章会讲塑料这类高分子）。把这样一座“分子积木仓库”整桶整桶烧掉取暖，在化学家眼里多少有点奢侈——这也是能源转型的另一个理由：把燃烧让给阳光和风，把碳链留给材料。

\section{一座靠沸点分拣的塔}

地底下抽出来的原油，是一锅黑乎乎、黏稠稠的“大杂烩”：从 1 个碳到 40 个碳以上的烷烃、环状烃、芳香烃全混在一起，谁也直接用不了。炼油厂做的第一件大事，不是“制造”汽油，而是\textbf{分拣}——把混合物按沸点拆开。这套装置叫分馏塔。

原理朴素得像蒸包子。第 22 章说过，液体沸腾就是分子挣脱邻居的握手逃进空气；链越短握手越松，沸点越低。把原油加热到三百多摄氏度送进修辞上叫“塔”、实际是一根几十米高的钢柱里，底部热、顶部冷。蒸汽往上升，升到哪一层冷到它的沸点以下，就在哪一层凝回液体，被托盘接住：最耐热的沥青留在塔底，柴油、煤油在中间各层，汽油蒸气爬到高处，最轻的甲烷、乙烷从塔顶以气体形式离开。一座塔，把一锅杂烩分成了十几层“货架”。

\begin{qiaoliang}
为什么“沸点不同”就能分开，而不是混在一起永远不分家？关键在于蒸发和凝结是双向的：任何温度下，分子都在液面进进出出，沸点低的成分在气相里占的比例总比液相里高（第 31 章的平衡思想在这里一样管用）。分馏塔里每一层托盘都相当于让蒸汽重新凝结、再蒸发一次，几十层托盘就是几十次“重新分配”，低沸点成分一层层被“筛”向塔顶。层数越多，分得越纯。整个过程没有断开任何一根化学键——进去是哪些分子，出来还是哪些分子，只是按沸点站好了队。\textbf{分馏是物理变化，不是制造新物质。}
\end{qiaoliang}

但分拣有个现实问题：市场对各层货架的需求不均匀。原油里天然的长链偏多，而全世界汽车要喝的汽油（5～12 个碳）远远不够分。怎么办？炼油厂的第二件事才是“制造”：\textbf{裂化}——把卖不动的长链分子拆短；\textbf{重整}——把拆好的短链重新拗成更好用的形状。

裂化的做法是加热、加压，再请出催化剂（第 30 章讲过，催化剂只改变反应路径和速度，不改变化学平衡的位置），把一根 16 个碳的柴油分子在某个碳碳单键处“掰断”，变成比如一个 8 碳的烷烃加一个 8 碳的烯烃。这一步是真正的化学变化：键断了，新分子出现了。

重整解决的是另一个毛病——\textbf{爆震}。汽油发动机靠火花塞点火，火焰应该平稳地扫过气缸。但如果汽油里直链烷烃太多，混合气会在火焰到达之前被压缩得自己提前着火，“咣”地敲一下气缸，这叫爆震，既费油又伤发动机。人们发现：同样是 8 个碳，直直一条链的正辛烷特别容易爆震，而拧出好几个支链的异辛烷非常抗爆。于是炼油厂用重整工艺把直链拗成支链和环状，汽油的抗爆性就上去了。

工程上给这个性能定了一把尺子，叫\textbf{辛烷值}：规定异辛烷的抗爆性为 100，正庚烷为 0，待测汽油跟这两种参照物的混合液比一比，抗爆性相当于多少比例的异辛烷，辛烷值就是多少。加油站标的 92 号、95 号，指的就是这把尺子上的读数——不是“纯度 95\%”，也不是“能量多 3 个单位”，只是抗爆性的排位。标号更高的汽油不一定让你的车更有劲，只是更不容易在高压发动机里提前自燃。

\section{燃烧：一场能量结算}

现在轮到本章的主角反应登场。天然气烧起来，化学方程式只有一行：

\[\chem{CH_4} + \chem{2O_2} \rightarrow \chem{CO_2} + \chem{2H_2O}\]

丁烷完全燃烧也一模一样，只是系数变大：

\[\chem{2C_4H_{10}} + \chem{13O_2} \rightarrow \chem{8CO_2} + \chem{10H_2O}\]

看产物：只有二氧化碳和水。这就是\textbf{完全燃烧}——氧气充足时，每个碳都抓到两个氧，每个氢都抓到自己的氧原子组成水，烧得干干净净。燃气灶的蓝火苗、盘底的水珠，都是它在场。

能量从哪来？第 27 章已经算过这笔账，这里再强调一次，因为它是最容易记错的地方：\textbf{断键要花钱，成键才进账}。点燃甲烷，先要吸热拆断 \chem{C-H} 键和 \chem{O=O} 键——所以要用火柴或电火花“点燃”这把火（第 29 章的活化能）；然后碳和氧、氢和氧结合成 \chem{CO_2} 和 \chem{H_2O}，新键形成时放出热。进账远大于支出，净结果就是放热，而且放出的热又替旁边的分子付了“点火费”，火就自己烧下去了。火焰的能量不是藏在燃料里“存着的热”，而是\textbf{产物的化学键比反应物的键更省钱}，差价以光和热的形式找零给你。

可要是氧气不够呢？碳抢不到足够的氧，只能烧到一半。产物里就会出现两种“半成品”：一氧化碳 \chem{CO}——碳只抓到一个氧；以及干脆没烧的碳颗粒，也就是你勺子上那层黑灰。这叫\textbf{不完全燃烧}。蜡烛的黄色火苗、熏黑的锅底、柴油机的黑烟，全都是它。

那抹黄色其实是一场意外的“灯光秀”：碳颗粒在火焰里被烧到上千摄氏度，像烧红的铁块一样发光（第 8 章讲过，足够热的物体都会发光），无数微小碳粒同时发光，火焰就成了明亮的黄色。蓝火苗里没有碳粒发光，只有气体反应时微弱的分子发光，所以是安静的蓝。开头的第一个问题，答案已经露头了。

不完全燃烧不只是浪费燃料。那两种半成品里，无色无味的一氧化碳是真正的杀手，值得单独说。

\section{一氧化碳：堵在车站的乘客}

你身体里的氧气运输系统是一套精密的物流：肺是装货站，血液里数以亿计的血红蛋白是货车，每辆货车中央有一个铁原子做的“车位”，恰好能抱住一个氧分子 \chem{O_2}（第 15 章讲配位化学时细说过这个结构），运到全身各处的组织卸货，再空车返回。

麻烦在于，一氧化碳分子的形状和电荷分布跟氧分子很像，也能卡进同一个车位——而且卡得比氧分子牢两百多倍。吸入一氧化碳后，一辆辆货车被它长期霸占，装不上氧，血液这条运输线就瘫痪了。大脑和心脏最先用氧量告急：先是头痛、头晕、恶心，接着失去意识，严重时几十分钟内致命。最阴险的是，人在睡梦中中毒时连“头痛”这个警报都收不到。

\begin{wuqu}
“燃气有味，漏了我能闻出来；一氧化碳中毒也能提前察觉。”——错得危险。天然气和液化石油气的主要成分甲烷、丙烷、丁烷，本身\textbf{都没有气味}；一氧化碳同样无色无味。你在家里闻到的“煤气味”，是燃气公司\textbf{故意加进去的}微量含硫臭味剂（硫醇类），臭到极低的浓度人也能察觉——这是人为安装的“嗅觉报警器”。但燃烧不充分产生的一氧化碳没有任何臭味剂保护，你闻不到它。真正可靠的防线只有三条：保持通风、不要在密闭房间里用燃气热水器或烧炭取暖、装一只一氧化碳报警器。冬天在门窗紧闭的屋里烧炭火锅，就是在跟这套运输系统赌博。
\end{wuqu}

所以请记住一条铁律：\textbf{任何火焰都不能待在密闭房间里}。燃气热水器必须接排风管，炭炉烧烤必须在室外，汽车不能在关着门的车库里怠速。这些事本章不给“在家试试”的环节——一氧化碳的危险实验只许看视频或老师演示。

\section{石油从哪里来}

\begin{zhengju}
石油埋在地下几百米到几千米的岩石里，它的“出生证明”是怎么被读出来的？总不能挖个洞下去看。

20 世纪 30 年代，德国化学家阿尔弗雷德·特赖布斯（Alfred Treibs）做了一件像侦探做的事：他从石油和油页岩里分离出了一类分子——卟啉。卟啉这个名字你不必记，但它是生命体里两种关键分子的骨架：叶绿素（植物抓住阳光的分子）和血红素（血红蛋白里含铁的那部分，前一节刚见过）都是卟啉家族成员。特赖布斯发现，石油里的卟啉结构和叶绿素、血红素的骨架几乎一一对应，只是中心的金属原子换了、侧链被岁月修剪过。这就像在事故现场找到了遇难者的身份证碎片：石油里躺着叶绿素和血红素的“化石”。

反对的声音当然也有：会不会是石油恰好合成了相似的结构？后来证据越堆越多。人们在石油里找到成百上千种“生物标志物”——只有生物才会制造的分子化石，比如一类叫藿烷的分子，只来自细菌的细胞膜。碳同位素也投了票：第 7 章讲过同位素，碳有碳-12 和碳-13 两种常见形式，生物体内的酶干活时偏爱轻的碳-12，所以生物来源的碳总是“偏轻”；石油的碳同位素比例恰好带着这种生物指纹，和无机岩石里的碳明显不同。

到今天，证据链闭合了：石油主要是古代海洋里浮游生物和藻类的遗体，被沉积物掩埋后，在缺氧、高温、高压下经过几百万到几亿年缓慢改造而成；天然气常常伴生，是这场改造中分解出的最轻部分。煤炭则主要来自古代陆地植物。所以它们被合称为\textbf{化石燃料}——字面上的“化石做的燃料”。你点燃的每一升汽油，都是几亿年前某个晴朗日子里被浮游生物存起来的阳光。
\end{zhengju}

\section{烧掉的代价：碳的搬家与温室效应}

化石燃料燃烧本质上是把埋了几亿年的碳，在几百年里集中搬回大气。来做一个真实的数量级估算：一辆家用车一年大约烧掉 1000 升汽油。汽油按辛烷 \chem{C_8H_{18}} 估算，密度约每升 0.75 千克，1000 升就是 750 千克。每个辛烷分子含 8 个碳，燃烧后变成 8 个 \chem{CO_2}。碳的原子量是 12，\chem{CO_2} 的式量是 44，也就是每 12 份碳变成 44 份二氧化碳。750 千克汽油里约八成五是碳，约 640 千克；乘以 $44/12$，得到约 \textbf{2300 千克二氧化碳}——这辆车一年往大气里倒了 2.3 吨看不见的气体，比车自身还重。全球十几亿辆车，再加上电厂、轮船、飞机，每年是人类活动往大气里搬的几百亿吨碳。

这些二氧化碳去了哪、干了什么？第 55 章会完整讲碳循环，这里先记住图景：碳在大气、海洋、陆地生物圈之间不停循环，燃烧化石燃料是在循环之外\textbf{额外加塞}。二氧化碳有个关键性质：它放行阳光里的可见光，却吸收地面向外辐射的红外线，像给地球多盖了一层被子——这就是\textbf{温室效应}。要强调两点边界：第一，温室效应本身是好事，没有它地球平均温度会是零下十几度；问题出在“加得太快”，大气 \chem{CO_2} 浓度从工业革命前的约 280 ppm 升到今天超过 420 ppm，被子加厚得太急，气候系统来不及平滑适应。第二，气候系统极其复杂，云、洋流、冰雪反馈交织，具体某年某地会怎样，科学只能给概率范围——但“人类烧化石燃料正在推高全球平均温度”这个主结论，证据等级和“吸烟有害健康”相当。

出路不是“停止用能量”，而是换碳的来源和去处：风电、光伏、核电不搬碳；电动车把排放从千万个排气管集中到可控的电厂；把捕集的 \chem{CO_2} 重新压回地下，或者用绿氢和 \chem{CO_2} 重新合成燃料，让碳循环闭合。这些\textbf{能源转型}技术每一样都建立在本书前面讲过的原理上——电池（第 34 章）、催化（第 30 章）、能量账（第 27 章）。化学没有制造这个问题，但解决这个问题绕不开化学。

\section{边界：本章模型在哪里会失灵}

趁还记得，给本章的模型划几条边界：

\begin{itemize}[itemsep=2pt]
\item \textbf{“烷烃=整条链按沸点排队”是简化。}真实的汽油是上百种分子的混合物，支链异构体的沸点和直链不同，所以表中每个馏分都给的是沸点范围而不是一个数。
\item \textbf{辛烷值是经验标度，不是分子性质的精确度量。}它只在标准发动机测试条件下有意义，乙醇等含氧燃料可以“超纲”地超过 100；它也不代表能量高低——高标号汽油不更“有劲”。
\item \textbf{完全燃烧是理想极限。}真实火焰里完全与不完全燃烧总是同时发生，只是比例不同；汽车尾气里永远有一氧化碳和碳颗粒，所以才有三元催化器这个“尾气化工厂”（催化剂，第 30 章）。
\item \textbf{“温室效应=烧燃料的全部后果”不对。}燃烧还带来颗粒物、氮氧化物、硫氧化物等近处污染，它们影响的是你的肺而不是气候，治理逻辑也不同。
\end{itemize}

\section{回到厨房那两团火}

现在可以兑现开场承诺了。

\textbf{蓝火苗和黄火苗：}燃气灶把天然气和空气预先混合好，氧气充足，甲烷烧得干净，没有碳颗粒发光，只剩下分子反应的淡蓝；蜡烛的蜡蒸气在焰心里缺氧，碳先聚成微粒，微粒被烧热发出黄光，飘到火焰外缘才烧尽——所以蜡烛是“一半照明、一半燃烧”的火。开头说的打火机黄火苗同理：丁烷直接喷出，来不及混够空气。

\textbf{打火机里的液体为什么一出来就成气：}丁烷沸点约 $-1\,^{\circ}\mathrm{C}$，在常温常压下本该是气体；打火机里加压把它摁成了液体（第 22 章，压强帮助分子间握手），一按阀门压强消失，它立刻沸腾成气喷出。那罐“像水的液体”，是被压强暂时囚禁的气体。

\textbf{燃气的臭味：}天然气本身无味，臭味是人为添加的硫醇——燃气公司替你装上的化学警报器。而燃烧不全产生的一氧化碳没有任何气味，只能靠通风和报警器。

三团火、三个问题，背后是同一套逻辑：\textbf{分子的结构决定它怎么烧，烧得干不干净决定它留下什么}。

\section{燃料之外：下一个接口}

这一章我们认识了几乎没有反应接口的烷烃——它们把全身“接口预算”都省下来，只干燃烧这一件事。可只要在碳链末端装上一个羟基（\chem{-OH}，第 37 章介绍过的官能团），分子的性格就全变了：它能溶于水，能散发香气，能被你肝脏里的酶认出来、拆掉，也能让微生物为之疯狂——这个叫羟基的小接口，制造出了酒、醋、香水和消毒水。下一章，我们给碳链接上第一个官能团，看看接口怎样改变一切。

\begin{tiaozhan}
用键能亲手结一次燃烧的账（跳过不影响主线）。查表可得近似键能（单位 $\mathrm{kJ/mol}$）：\chem{C-H} 约 413，\chem{O=O} 约 498，\chem{C=O}（在 \chem{CO_2} 中）约 799，\chem{O-H} 约 463。甲烷燃烧 $\chem{CH_4} + \chem{2O_2} \rightarrow \chem{CO_2} + \chem{2H_2O}$：断键支出为 $4\times 413 + 2\times 498 = 2648$；成键收入为 $2\times 799 + 4\times 463 = 3450$。净放热约 $3450-2648=802$ $\mathrm{kJ/mol}$，与实测值 890 相当接近（差距来自键能是平均值的近似）。注意方向：\textbf{支出在前、收入在后}——所以火焰必须先被点燃（付活化能），然后才自己滚起来。再想一想：为什么一氧化碳还能继续燃烧放热？因为 \chem{CO} 到 \chem{CO_2} 还能再成一根新键，这正是不完全燃烧“浪费掉”的那部分能量。
\end{tiaozhan}

\section*{练一练}

\begin{sikaoti}
\item 丙烷 \chem{C_3H_8}（液化石油气的主要成分）完全燃烧的方程式是什么？配平后数一数：每烧 1 个丙烷分子要消耗几个氧分子？如果氧气减半，碳元素可能变成哪两种产物？
\item 画一幅“分馏塔示意图”：画一根竖塔，标出底部进料口和从上到下的温度变化，再把天然气、汽油、煤油、柴油、沥青五类产物画在各自出料的高度上。图旁注明：分拣靠的是哪个性质的差异？
\item 汽油标号从 92 换成 95，汽车并不会更有劲。用本章的模型解释：辛烷值量的到底是什么？什么结构的分子辛烷值高？
\item 参考本章的数量级估算，算一算：你家若一个月用 20 立方米天然气（按纯甲烷计，每立方米约 0.72 千克），完全燃烧后产生多少千克二氧化碳？（提示：\chem{CH_4} 式量 16，其中碳占 12。）
\item 冬天有人在密闭房间里烧炭取暖，说“我开着一条门缝，而且没闻到什么味道，应该没事”。请用粒子层的语言，逐条指出这句话里的两处错误。
\item 石油卟啉证据故事里说“生物的碳偏轻”。如果未来在火星岩石里发现某种含碳分子，它的碳同位素比例明显偏轻，这能证明火星有生命吗？还需要排除哪些替代解释？
\end{sikaoti}
